\documentclass[11pt]{article}
\usepackage[a4paper,margin=26mm]{geometry}
\usepackage{amsmath,amssymb,mathtools}
\usepackage{unicode-math}
\setmainfont{Latin Modern Roman}
\setmathfont{Latin Modern Math}
\usepackage{xurl,hyperref}
\hypersetup{hidelinks,pdftitle={Problem statement}}
\newcommand{\sref}[2]{\hyperlink{source:#1:#2}{[S#2]}}
\newcommand{\cataloguescope}[1]{#1}
\begin{document}
\section{Branched surfaces and triangulations — Question 7.8}\label{statement}

\subsection{Definitions and mathematical statement}

\paragraph{Shared notation.}$\mathbb N=\{1,2,\ldots\}$, and $\mathbb Z,\mathbb Q,\mathbb R,\mathbb C$ denote the integers, rationals, reals and complex numbers. Any other conventions are specified locally.


A branched surface $B\subset M$ is a smoothly tangent $2$-complex with the standard local branching models. Its fibered neighborhood $N(B)$ is an interval bundle with the usual identifications over the branch locus. A lamination is carried by $B$ if it lies in $N(B)$ transverse to interval fibers, and fully carried if it meets every fiber. Splitting means cutting this fibered neighborhood along carried surface pieces and taking the induced branched surface. The branch locus has a direction in which branches merge. A sink disk is an embedded disk in the branched surface whose boundary is in that locus and whose branch directions all point inward. An automatic lamination consists of a fixed cell decomposition of $B$, a lamination $\Lambda$ fully carried by $B$, and a finite-state automaton recognizing combinatorial paths in $B$ that represent geodesic paths in its leaves.
\paragraph{Statement recorded in the supplied source.}
Does every branched surface without a sink disk carry an automatic lamination in this sense? In particular, the finite-state recognition is required, beyond merely carrying a lamination. \sref{statement}{2}
\subsection{Short English statement}
Does absence of sink disks ensure that a branched surface carries a lamination whose leafwise geodesic paths are recognized by a finite-state automaton?
\subsection{Sources}
\begin{description}
\item[\hypertarget{source:statement:1}{S1}] ULAM.AI and the UnsolvedMath Contributors, \emph{UnsolvedMath: A Curated Collection of Open Mathematics Problems}, catalogue entry. CC BY 4.0. \url{https://huggingface.co/datasets/ulamai/UnsolvedMath}.
\item[\hypertarget{source:statement:2}{S2}] Danny Calegari, \emph{Problems in foliations and laminations of 3-manifolds}, version 0.78 (2002), Question 7.8; Definitions 1.1--1.10 and local definitions and remarks. \url{https://arxiv.org/abs/math/0209081v1}.
\end{description}
\label{end:statement}
\end{document}
